% DRAFT — placeholder for post-experiment fill.
% Will be spliced into 4_experiments.tex after \subsection{Cross-benchmark generalisation: MMAD}.

\subsection{A unified task-aware agent for MMAD}
\label{ssec:mmad_fulltype}

The cross-benchmark MMAD evaluation in \S\ref{ssec:mmad} used only the
Anomaly-Detection subset (Yes/No). MMAD additionally contains eight
$4$-way multiple-choice question types that probe \emph{defect
classification} (\num{4688} questions), \emph{defect localization}
(\num{4878}), \emph{defect description} (\num{4710}), \emph{defect
analysis} (\num{4782}), and four \emph{object}-understanding types
(\num{12315} combined). We extend our agent to a \emph{unified
task-aware} output head that accepts an optional question and a set of
multiple-choice options and produces, in a single forward pass, (i) a
continuous anomaly score (for ensembling on AD), (ii) a letter answer
for MCQ, and (iii) free-text when no options are present. The mode is
inferred from the input inside the first cognitive turn---no separate
routing call---so the same agent that produces anomaly scores on
CrossDomainVAD also answers MCQ items on MMAD.

\paragraph{Setup.}
We fit per-class decision thresholds for the AD subset on a
stratified $n{=}\TODO_DEV$-image \emph{dev} sample (all 38 product
classes, seed $= 42$), then freeze and evaluate on a held-out
$n{=}\TODO_TEST$-image stratified \emph{test} sample (different seed).
Ensembling is applied only to AD questions; for the remaining eight
MCQ types we compare the unified agent's letter prediction against
the same Direct~VLM baseline with a strictly comparable MCQ-aware
prompt. Backbone: Qwen3.5-VL-27B, temperature~$0$, same $4$ normal
reference templates as \S\ref{ssec:mmad}.

\begin{table}[h]
\centering
\caption{\textbf{MMAD full-type MCQ accuracy} (Qwen3.5-VL-27B, test
sample; $n$ per type in parentheses). For Anomaly~Detection we
additionally report the $0.5$-averaged ensemble and its per-class
calibrated variant. For the eight non-AD MCQ types we report the
unified agent alone (Direct~VLM baseline next to it).}
\label{tab:mmad_fulltype}
\small
\begin{tabular}{lcccc}
\toprule
Question type ($n$)            & Direct VLM & Unified Agent the refutation agent & Ensemble & Ens.~(per-class $\tau$) \\
\midrule
Anomaly Detection (\TODO)      & \TODO & \TODO & \TODO & \TODO \\
\midrule
Defect Classification (\TODO)  & \TODO & \TODO & --- & --- \\
Defect Localization (\TODO)    & \TODO & \TODO & --- & --- \\
Defect Description (\TODO)     & \TODO & \TODO & --- & --- \\
Defect Analysis (\TODO)        & \TODO & \TODO & --- & --- \\
\midrule
Object Classification (\TODO)  & \TODO & \TODO & --- & --- \\
Object Analysis (\TODO)        & \TODO & \TODO & --- & --- \\
Object Structure (\TODO)       & \TODO & \TODO & --- & --- \\
Object Details (\TODO)         & \TODO & \TODO & --- & --- \\
\midrule
Pooled (\TODO total)           & \TODO & \TODO & --- & --- \\
\bottomrule
\end{tabular}
\end{table}

\paragraph{Findings.}
% Fill after experiments:
% 1. Unified agent closes the MCQ-accuracy gap on AD (from $-0.8$ pp to
%    $\Delta{>}0$) once per-class thresholds are calibrated, while
%    retaining the AUROC gain of \S\ref{ssec:mmad}.
% 2. On Defect-* types the agent's per-image tool reasoning is
%    informative and delivers $+\TODO$ pp over Direct.
% 3. On Object-* types the agent matches Direct within noise, confirming
%    that the mode-inference step does not sacrifice non-detection
%    performance.
%
% If positive numbers are obtained, position this as evidence that a
% single trained-free agent can span detection and fine-grained MCQ
% answering without per-task retraining.
